% ======================================================================
% Appendix L. Trace class criteria and regularization
% PATH: src/appendices/L-trace-class-criteria-and-regularization.tex
% ======================================================================

\appendix
\chapter{Trace class criteria and regularization}
\label{app:L.traceclass}

\section{Purpose and scope}
\label{sec:L.purpose}

This appendix collects functional-analytic tools used throughout the monograph to justify trace identities for
localized spectral projectors and for their microlocal perturbations.  The principal operators in the main text are
constructed by applying smooth functional calculus to $\sqrt{\Delta}$ and by inserting semiclassical cutoffs and
pseudodifferential weights.  In noncompact finite-volume settings, traces must be interpreted in a renormalized
sense due to cusp divergences, and the continuous spectrum requires a canonical subtraction procedure.

The purpose of this appendix is to provide a unified set of criteria ensuring that the operators encountered in
Chapters~\ref{chap:05-parametrix}--\ref{chap:11-final} are of trace class (or admit a well-defined regularized trace)
and that the regularization is stable under the admissible perturbations used in the localized trace formula.

The statements are standard but are recorded here in a form adapted to the present framework.

\section{Operator ideals and basic trace identities}
\label{sec:L.ideals}

\subsection{Schatten classes}
\label{subsec:L.schatten}

Let $\mathcal{H}$ be a complex separable Hilbert space and let $\mathcal{B}(\mathcal{H})$ denote the bounded
operators on $\mathcal{H}$.  For $1\le p<\infty$, the Schatten class $\mathcal{S}_p(\mathcal{H})$ consists of all
compact operators $T$ whose singular values $(s_n(T))_{n\ge 1}$ satisfy
\[
\|T\|_{\mathcal{S}_p}^p := \sum_{n\ge 1} s_n(T)^p < \infty.
\]
The trace class is $\mathcal{S}_1(\mathcal{H})$ and the Hilbert--Schmidt class is $\mathcal{S}_2(\mathcal{H})$.

If $T\in \mathcal{S}_1(\mathcal{H})$ the trace is defined by
\[
\Tr(T) := \sum_{n\ge 1} \langle Te_n,e_n\rangle,
\]
where $(e_n)$ is any orthonormal basis; the sum converges absolutely and is basis-independent.

\subsection{Cyclic property}
\label{subsec:L.cyclic}

We record the standard cyclic property of the trace in the only form we use.

\begin{lemma}[Cyclicity in Schatten classes]
\label{lem:L.cyclic}
Let $A\in \mathcal{B}(\mathcal{H})$ and $B\in \mathcal{S}_1(\mathcal{H})$. Then $AB,BA\in \mathcal{S}_1(\mathcal{H})$
and
\[
\Tr(AB)=\Tr(BA).
\]
More generally, if $A\in\mathcal{S}_p$, $B\in\mathcal{S}_q$ with $1/p+1/q=1$, then $AB\in\mathcal{S}_1$ and
$\Tr(AB)=\Tr(BA)$.
\end{lemma}

\begin{proof}
Standard; see any functional analysis reference on Schatten ideals.
\end{proof}

\subsection{Trace norm bounds}
\label{subsec:L.tracebounds}

\begin{lemma}[Ideal bounds]
\label{lem:L.idealbounds}
If $A\in\mathcal{B}(\mathcal{H})$ and $B\in\mathcal{S}_1(\mathcal{H})$, then
\[
\|AB\|_{\mathcal{S}_1} \le \|A\|_{\mathcal{B}(\mathcal{H})}\,\|B\|_{\mathcal{S}_1},
\qquad
\|BA\|_{\mathcal{S}_1} \le \|A\|_{\mathcal{B}(\mathcal{H})}\,\|B\|_{\mathcal{S}_1}.
\]
If $A\in\mathcal{S}_2$ and $B\in\mathcal{S}_2$, then $AB\in\mathcal{S}_1$ and
\[
\|AB\|_{\mathcal{S}_1}\le \|A\|_{\mathcal{S}_2}\,\|B\|_{\mathcal{S}_2}.
\]
\end{lemma}

\begin{proof}
Standard Hölder inequalities for Schatten norms.
\end{proof}

\section{Integral kernel criteria}
\label{sec:L.kernel}

\subsection{Hilbert--Schmidt and trace class kernels}
\label{subsec:L.hs.trace}

Let $X$ be a $\sigma$-finite measure space with measure $d\mu$ and let $\mathcal{H}=L^2(X,d\mu)$.
An operator $T$ is called an \emph{integral operator} if
\[
(Tf)(x)=\int_X K_T(x,y)f(y)\,d\mu(y)
\]
for some measurable kernel $K_T$.

\begin{lemma}[Hilbert--Schmidt criterion]
\label{lem:L.hs.criterion}
If $K_T\in L^2(X\times X)$, then $T\in \mathcal{S}_2(\mathcal{H})$ and
\[
\|T\|_{\mathcal{S}_2}^2 = \int_{X\times X} |K_T(x,y)|^2\,d\mu(x)\,d\mu(y).
\]
\end{lemma}

\begin{proof}
Standard.
\end{proof}

\begin{lemma}[Trace class via factorization]
\label{lem:L.trace.factor}
If $T=T_1T_2$ with $T_1,T_2\in\mathcal{S}_2(\mathcal{H})$, then $T\in\mathcal{S}_1(\mathcal{H})$ and
\[
\|T\|_{\mathcal{S}_1}\le \|T_1\|_{\mathcal{S}_2}\,\|T_2\|_{\mathcal{S}_2}.
\]
\end{lemma}

\begin{proof}
Immediate from Lemma~\ref{lem:L.idealbounds}.
\end{proof}

\subsection{Trace by diagonal restriction}
\label{subsec:L.diagonal}

In compact settings, the trace of an integral operator can be computed by integrating the diagonal of its kernel.
We record a safe version sufficient for the main text.

\begin{proposition}[Trace by diagonal integral]
\label{prop:L.diagonal.trace}
Let $X$ be a compact manifold with smooth density $d\mu$.
Suppose $T$ is a trace class operator on $L^2(X)$ with continuous kernel $K_T(x,y)$.
Then
\[
\Tr(T)=\int_X K_T(x,x)\,d\mu(x).
\]
\end{proposition}

\begin{proof}
Standard approximation by finite-rank operators.
\end{proof}

\begin{remark}
\label{rem:L.noncompact.diagonal}
In noncompact finite-volume settings, $K_T(x,x)$ may fail to be integrable due to cusp divergences even when $T$ is
trace class after renormalization.  The renormalized trace is addressed in
Section~\ref{sec:L.renormalized}.
\end{remark}

\section{Functional calculus and smoothing operators}
\label{sec:L.functional}

\subsection{Spectral calculus for $\sqrt{\Delta}$}
\label{subsec:L.spectralcalc}

Let $\Delta$ be the nonnegative Laplace--Beltrami operator on $L^2(X)$, with $X$ a finite-area hyperbolic surface.
For a bounded Borel function $f:\R_{\ge 0}\to\C$, the operator $f(\sqrt{\Delta})$ is defined via the spectral theorem.

In the compact case, $f(\sqrt{\Delta})$ is trace class whenever $f(\lambda)$ decays sufficiently fast.
In the finite-volume case, the operator is not trace class in general because of the continuous spectrum and the cusp
divergence of the heat kernel.  However, localized functional calculus yields trace class after applying either
compactly supported cutoffs or renormalized traces.

\subsection{Schwartz functional calculus}
\label{subsec:L.schwartz}

The principal operators of the localized trace method have the form
\[
T_{\lambda,\eta} := \chi\!\left(\frac{\sqrt{\Delta}-\lambda}{\eta}\right),
\]
where $\chi\in\mathcal{S}(\R)$ is even and $\eta>0$.

\begin{lemma}[Kernel smoothing]
\label{lem:L.kernel.smoothing}
Let $\chi\in\mathcal{S}(\R)$ and define $T_{\lambda,\eta}$ as above. Then $T_{\lambda,\eta}$ is a smoothing operator.
In particular, its Schwartz kernel $K_{\lambda,\eta}(x,y)$ is $C^\infty$ on $X\times X$.
\end{lemma}

\begin{proof}
This follows from the Fourier inversion formula
\[
\chi\!\left(\frac{\sqrt{\Delta}-\lambda}{\eta}\right)
=
\frac{1}{2\pi}
\int_{\R}
\widehat{\chi}(t)\,
e^{it(\sqrt{\Delta}-\lambda)/\eta}\,dt,
\]
combined with standard propagation properties of the wave group.
\end{proof}

\section{Pseudodifferential operators and Schatten bounds}
\label{sec:L.psiDO}

\subsection{Compact case: order bounds}
\label{subsec:L.compact.psiDO}

Let $M$ be a compact smooth manifold of dimension $d$ and let $\Psi^m(M)$ denote classical pseudodifferential
operators of order $m$.

\begin{theorem}[Schatten criterion for pseudodifferential operators]
\label{thm:L.schatten.psiDO}
If $A\in \Psi^m(M)$ with $m<-d/p$, then $A\in \mathcal{S}_p(L^2(M))$.
In particular, if $m<-d$ then $A$ is trace class.
\end{theorem}

\begin{remark}
\label{rem:L.psiDO.references}
This is a classical consequence of Weyl asymptotics for singular values and Sobolev embedding; see standard
treatments of pseudodifferential operator ideals.
\end{remark}

\subsection{Localized operators}
\label{subsec:L.localized}

The operators used in the localized trace formula are not simply pseudodifferential of negative order: they are
constructed by composing a functional calculus operator with microlocal cutoffs.  A typical example is
\[
A_{\lambda,\eta}(a)
=
\chi\!\left(\frac{\sqrt{\Delta}-\lambda}{\eta}\right)\Op(a),
\]
where $a\in C_c^\infty(T^*X)$ is a symbol supported away from the cusp region.

The relevant feature is that $\chi((\sqrt{\Delta}-\lambda)/\eta)$ is smoothing and $\Op(a)$ is bounded, hence
$A_{\lambda,\eta}(a)$ is smoothing.  In the compact case this is automatically trace class.

In finite-volume settings the issue is not smoothness but integrability of the diagonal of the kernel near the cusps,
which motivates renormalization.

\section{Renormalized traces on finite-volume surfaces}
\label{sec:L.renormalized}

\subsection{Cusp divergence and canonical subtraction}
\label{subsec:L.cuspdiv}

Let $X=\Gamma\backslash\mathbb{H}$ have cusps.
For smoothing operators $T$ with kernel $K_T(x,y)$, the function $x\mapsto K_T(x,x)$ may fail to be integrable due to
the asymptotic volume growth in cusp coordinates.  A renormalized trace is defined by truncation and subtraction of
explicit divergent terms.

Let $X(Y)$ be the truncated surface obtained by removing cusp neighborhoods above height $Y$.
Define
\[
\Tr_Y(T):=\int_{X(Y)} K_T(x,x)\,d\mu(x).
\]
Typically $\Tr_Y(T)$ diverges as $Y\to\infty$ like $c(T)\log Y$ plus bounded terms.

\begin{definition}[Renormalized trace]
\label{def:L.renorm.trace}
Suppose there exists a constant $c(T)$ such that
\[
\Tr_Y(T) = c(T)\log Y + \Tr^{\mathrm{ren}}(T) + o(1)
\qquad (Y\to\infty).
\]
Then $\Tr^{\mathrm{ren}}(T)$ is called the \emph{renormalized trace} of $T$.
\end{definition}

\begin{remark}
\label{rem:L.renorm.canonical}
In the present framework the subtraction constant $c(T)$ is not chosen ad hoc.  It is determined canonically by the
Maass--Selberg renormalization identity recorded in Appendix~\ref{app:K.scattering}.
\end{remark}

\subsection{Trace class vs.\ renormalized trace}
\label{subsec:L.traceclass.vs.renorm}

We emphasize that smoothing does not imply trace class in the finite-volume case.

\begin{proposition}[Smoothing does not imply trace class on cusped surfaces]
\label{prop:L.smoothing.not.trace}
Let $X$ be a finite-area hyperbolic surface with at least one cusp.
Then there exist smoothing operators $T$ on $L^2(X)$ such that $T\notin\mathcal{S}_1(L^2(X))$.
\end{proposition}

\begin{proof}[Proof sketch]
One may construct a smoothing operator whose diagonal kernel behaves like a nonintegrable constant in cusp
coordinates.  The integral over the cusp then diverges logarithmically.
\end{proof}

Thus, the correct statement in the cusped case is that the localized trace formula produces operators whose trace is
well-defined after renormalization, and that the renormalized trace coincides with the spectral trace expression
(discrete plus continuous spectrum) up to explicit cutoff-dependent constants.

\subsection{Compatibility with Maass--Selberg}
\label{subsec:L.ms.compat}

Let $T$ be an operator obtained by functional calculus applied to $\sqrt{\Delta}$ with test function $h$.
Then the continuous spectrum trace contribution is expressed through $\varphi'/\varphi$ as in
Appendix~\ref{app:K.scattering}.
The Maass--Selberg relations show that the truncation divergence is precisely the $\log Y$ term and that the
renormalized trace is independent of the truncation procedure.

\begin{theorem}[Canonical renormalized trace identity]
\label{thm:L.canonical.renorm}
Let $h$ be admissible and let $T=h(\sqrt{\Delta})$ be defined by spectral calculus.
Then the truncated trace $\Tr_Y(T)$ admits an expansion
\begin{equation}\label{eq:L.renorm.expansion}
\Tr_Y(T)
=
c_h\log Y
+
\Tr^{\mathrm{ren}}(T)
+
o(1),
\qquad Y\to\infty,
\end{equation}
where $c_h$ is explicitly determined by the scattering data and by the value $h(0)$.
Moreover, $\Tr^{\mathrm{ren}}(T)$ equals the spectral trace expression
\begin{equation}\label{eq:L.spectral.renorm}
\Tr^{\mathrm{ren}}(T)
=
\sum_{j} h(r_j)
+
\frac{1}{4\pi}
\int_{-\infty}^{\infty}
h(t)\,
\left(
-\frac{\varphi'}{\varphi}\!\left(\frac12+it\right)
\right)\,dt
+
\mathcal{R}_h,
\end{equation}
where $\mathcal{R}_h$ is an explicit cutoff correction determined by the normalization convention.
\end{theorem}

\begin{remark}
\label{rem:L.renorm.interface}
The term $\mathcal{R}_h$ is precisely the correction isolated in Chapter~\ref{chap:08-normalization}.
The monograph is structured so that all such corrections are consolidated into a fixed normalization operator
$h\mapsto h^\sharp$.
\end{remark}

\section{Regularization by analytic continuation}
\label{sec:L.zeta.reg}

\subsection{Zeta-regularized traces}
\label{subsec:L.zeta}

In some contexts it is convenient to define traces through analytic continuation of a zeta function.  Let $A$ be a
positive self-adjoint operator with discrete spectrum $(\lambda_n)$.
One defines
\[
\zeta_A(s)=\sum_n \lambda_n^{-s},
\]
for $\Re(s)$ large, and extends meromorphically.

The zeta-regularized trace of $f(A)$ is then extracted by Mellin inversion.  In the present work, however, the
localized trace formula does not rely on such a procedure: the regularization is performed geometrically by cusp
truncation and canonically by Maass--Selberg.

We record only the compatibility statement.

\begin{proposition}[Compatibility of renormalized and zeta regularizations]
\label{prop:L.compat.zeta}
In the compact case, the renormalized trace coincides with the ordinary trace and with the zeta-regularized trace.
In the finite-volume cusped case, the renormalized trace agrees with the analytic continuation regularization
whenever both are defined and normalized by the same scattering subtraction.
\end{proposition}

\section{Trace class stability under perturbations}
\label{sec:L.stability}

A crucial point in the main text is that trace identities remain valid after inserting microlocal weights and after
localizing in shrinking spectral windows.  This requires stability of the trace class (or renormalized trace class)
under perturbations.

\subsection{Perturbation lemma: bounded conjugation}
\label{subsec:L.conjugation}

\begin{lemma}[Bounded conjugation stability]
\label{lem:L.conjugation}
Let $T\in \mathcal{S}_1(\mathcal{H})$ and $B\in\mathcal{B}(\mathcal{H})$.
Then $BTB^{-1}\in\mathcal{S}_1(\mathcal{H})$ whenever $B$ is invertible, and
\[
\Tr(BTB^{-1})=\Tr(T).
\]
\end{lemma}

\begin{proof}
By cyclicity of trace.
\end{proof}

\subsection{Perturbation lemma: commutators}
\label{subsec:L.commutators}

The localized trace method repeatedly uses that commutator errors are trace class and have controlled trace norm.

\begin{lemma}[Commutator trace bound]
\label{lem:L.commutator}
Let $A\in\mathcal{S}_2(\mathcal{H})$ and $B\in\mathcal{S}_2(\mathcal{H})$. Then
\[
[A,B]=AB-BA\in\mathcal{S}_1(\mathcal{H})
\]
and
\[
\|[A,B]\|_{\mathcal{S}_1}
\le
2\|A\|_{\mathcal{S}_2}\|B\|_{\mathcal{S}_2}.
\]
\end{lemma}

\begin{proof}
Immediate from Lemma~\ref{lem:L.idealbounds}.
\end{proof}

\subsection{Stability of renormalized traces}
\label{subsec:L.renorm.stability}

\begin{proposition}[Renormalized trace stability under compactly supported perturbations]
\label{prop:L.renorm.stability}
Let $T$ be a smoothing operator on $X$ whose kernel is supported away from the cusps.
Then $T$ is trace class and its trace coincides with its renormalized trace.
More generally, if $S$ is any smoothing operator and $T$ is compactly supported away from cusps, then
\[
\Tr^{\mathrm{ren}}(S+T)=\Tr^{\mathrm{ren}}(S)+\Tr(T).
\]
\end{proposition}

\begin{proof}
Compact support removes the cusp divergence and reduces the renormalized trace to the usual diagonal integral.
\end{proof}

\section{Trace class criteria for the localized projector kernels}
\label{sec:L.projector.criteria}

We now record the exact criterion used implicitly in the localized trace formula.

\begin{theorem}[Trace class criterion for localized spectral projectors]
\label{thm:L.localized.traceclass}
Let $h\in\mathcal{S}(\R)$ be even and define
\[
T_h := h(\sqrt{\Delta}).
\]
Then:
\begin{enumerate}
\item If $X$ is compact, then $T_h$ is trace class and $\Tr(T_h)=\int_X K_h(x,x)\,d\mu(x)$.
\item If $X$ has cusps, then $T_h$ admits a canonical renormalized trace $\Tr^{\mathrm{ren}}(T_h)$, and the spectral
decomposition identity \eqref{eq:L.spectral.renorm} holds with continuous spectrum correction expressed by
$\varphi'/\varphi$.
\end{enumerate}
\end{theorem}

\begin{remark}
\label{rem:L.projector.use}
This theorem is the analytic backbone allowing the passage from wave kernel parametrix constructions to global trace
identities.  All subsequent localization steps amount to choosing $h$ of the form
$h(t)=\chi((t-\lambda)/\eta)$, for which the trace estimates are uniform in $(\lambda,\eta)$ in the admissible regime.
\end{remark}

\section{Regularization via smoothing removal}
\label{sec:L.smoothing.removal}

The final step in the main text removes smoothing cutoffs via Tauberian arguments.  From the perspective of trace
class theory, this corresponds to approximating discontinuous spectral projectors by Schwartz test functions.

We record the abstract approximation lemma used repeatedly.

\begin{lemma}[Approximation of spectral projectors]
\label{lem:L.approx.projectors}
Let $f\in L^\infty(\R)$ be bounded and compactly supported, and let $(f_n)\subset \mathcal{S}(\R)$ satisfy
$f_n\to f$ pointwise and $\sup_n\|f_n\|_\infty<\infty$.
Then $f_n(\sqrt{\Delta})$ converges to $f(\sqrt{\Delta})$ in the strong operator topology.
If, in addition, $X$ is compact and $f$ has sufficient decay, then the traces converge.
\end{lemma}

\begin{remark}
\label{rem:L.tauberian.bridge}
In the noncompact case, trace convergence must be interpreted at the level of renormalized traces, and the relevant
estimates are transferred to the Tauberian removal mechanism in Appendix~\ref{app:D.tauberian}.
\end{remark}

\section*{Bibliographic notes}

Trace ideals and Schatten class bounds are standard; see any modern text on operator theory.
The role of trace class criteria in spectral geometry is classical and is closely tied to heat kernel asymptotics.
For finite-volume hyperbolic surfaces, the necessity of renormalized traces is well-known and is canonically governed
by the Maass--Selberg relations and the scattering determinant, summarized in Appendix~\ref{app:K.scattering}.
For detailed treatments of renormalized trace formulae in hyperbolic settings, see \cite{Hejhal1983,Iwaniec2002}.

% ======================================================================
% End of L-trace-class-criteria-and-regularization.tex
% ======================================================================
